%% SinglePress — Lossless compression of sparse count matrices
%% using Value-Ordered Compressed Sparse Column encoding
\documentclass[10pt,twocolumn,letterpaper]{article}
\usepackage{singletai-preprint}
\usepackage{graphicx}
\usepackage{amsmath,amssymb}
\usepackage{booktabs}
\usepackage{natbib}
\usepackage{xcolor}
\usepackage{multirow}
\bibliographystyle{unsrtnat}

% ── Metadata ─────────────────────────────────────────────────────
\title{SinglePress: Lossless Compression of Sparse\\Count Matrices via Value-Ordered Column Encoding}
\author{
  Zach DeBruine\textsuperscript{1,*}
  \\[6pt]
  \textsuperscript{1}Singlet Bio
}

\begin{document}
\maketitle
\correspondingauthor{zach@singlet.ai}

\begin{abstract}
Sparse count matrices from single-cell RNA sequencing (scRNA-seq) are a dominant data type in modern genomics, with public repositories exceeding 500 million cells.
Standard CSC (Compressed Sparse Column) storage treats integer counts as 64-bit floats and encodes row indices as absolute 32-bit integers, ignoring the extreme value concentration and spatial structure inherent to count data.
We present VOCSC (Value-Ordered Compressed Sparse Column), a lossless compression scheme that exploits three statistical properties of scRNA-seq matrices: (1) extreme value concentration (72\% of nonzeros equal 1), (2) spatial locality of row positions within each value group, and (3) compressibility of gap sequences under byte-plane separation.
VOCSC groups column entries by value, encodes positions as gap-1 sequences, applies adaptive byte-plane splitting, and compresses with zstd at level~3.
On a corpus of 20 scRNA-seq count matrices, VOCSC achieves a median compression ratio of 12.5$\times$ over standard CSC (float64 values + int32 indices), decoding at 4,624~MB/s and encoding at 240~MB/s on 8 CPU threads.
This represents a 2.2$\times$ improvement over generic zstd applied directly to CSC arrays (5.85$\times$) and a 2.4$\times$ improvement over gzip (5.42$\times$).
We characterize the contribution of each pipeline stage through a systematic ablation study, validate behavior on synthetic matrices with controlled statistical properties, and develop a predictive compressibility model ($R^2 = 0.986$) relating matrix statistics to achievable compression ratio.
We further evaluate but ultimately reject higher-effort codec variants (zstd-22, Brotli-11) as yielding marginal ratio gains ($<$9\%) at 7--40$\times$ slower encoding and up to 3$\times$ slower decoding.
We also evaluate three extended encoding strategies---dense head rows, implicit value labels, and a hybrid scheme---finding consistent but modest improvements (median $-$1.3\% to $-$1.7\%) that do not justify the added codec complexity.
\end{abstract}

\keywords{lossless compression \and sparse matrices \and single-cell genomics \and CSC format \and data compression}

% ══════════════════════════════════════════════════════════════════
\section{Introduction}
% ══════════════════════════════════════════════════════════════════

Single-cell RNA sequencing (scRNA-seq) has become the workhorse assay of modern molecular biology, generating count matrices that record transcript abundance per gene per cell~\cite{zheng2017massively,svensson2018exponential}.
A typical experiment produces a matrix with $10^4$--$10^5$ genes (rows) and $10^3$--$10^6$ cells (columns), with 95--99.9\% of entries equal to zero~\cite{bacher2016design}.
Public repositories such as the Gene Expression Omnibus~\cite{edgar2002gene} and the Human Cell Atlas~\cite{regev2017human} now collectively exceed 500 million cells, creating an acute need for efficient storage and retrieval.
These matrices are universally stored in CSC (Compressed Sparse Column) format, where nonzero entries are represented as tuples of (row index, value) organized by column through an index pointer array.
The standard on-disk representation uses 32-bit integers for row indices and 64-bit floating-point for values, yielding a per-nonzero cost of $12$ bytes plus $4$ bytes per column for the pointer array.

Despite the structured statistical properties of count data---dominated by small integers, with strong spatial patterns---general-purpose compressors such as gzip~\cite{deutsch1996gzip}, LZ4~\cite{collet2013lz4}, and Zstandard (zstd)~\cite{collet2018zstandard} treat these byte streams as opaque.
Applied to raw CSC arrays, they achieve only modest compression (2.7--5.9$\times$) because the interleaved integer and floating-point bytes disrupt pattern detection.
Specialized formats such as HDF5/AnnData~\cite{anndata,folk2011overview} wrap CSC with library-level compression but do not exploit domain structure.

We observe three statistical regularities in scRNA-seq count matrices that a domain-informed codec can exploit:

\begin{enumerate}
\item \textbf{Value concentration.} On our 20-matrix corpus, a median of 76.5\% of nonzeros equal 1, and 92\% are $\leq 3$. Although matrices contain up to thousands of distinct values globally, individual columns---which are processed independently---exhibit far fewer distinct values due to the small counts typical of unique molecular identifier (UMI) data~\cite{islam2014quantitative}.
\item \textbf{Implicit values.} Since values are concentrated and few, grouping entries by value within each column eliminates the explicit value stream entirely. The value becomes implicit from group membership.
\item \textbf{Gap compressibility.} Within a value group, row positions form a monotonically increasing sequence well-approximated by a quasi-uniform distribution. The gap-1 transform (successive differences minus 1) produces low-entropy sequences amenable to byte-plane separation and entropy coding.
\end{enumerate}

Our codec, VOCSC, implements a four-stage pipeline: (1) row reordering by nonzero density, (2) value-ordered column encoding with gap-1 index representation, (3) adaptive byte-plane splitting (2 planes for 16-bit gaps, 4 for 32-bit), and (4) zstd compression at level 3.
We show that this combination achieves 12.5$\times$ compression over raw CSC while maintaining 4,624~MB/s decode throughput---suitable for interactive analysis of large-scale atlases.

Our contributions are:
\begin{itemize}
\item The VOCSC encoding scheme and its efficient C++/OpenMP implementation.
\item A systematic ablation quantifying the contribution of each pipeline stage.
\item Synthetic-data experiments confirming that compression ratio varies predictably with matrix statistics.
\item A linear predictive model ($R^2 = 0.986$) estimating achievable compression from matrix properties.
\item A comprehensive comparison against generic compressors and an analysis of higher-effort codec variants.
\item An evaluation of extended encoding strategies (dense head, implicit values, hybrid) that quantifies the diminishing returns of further structural exploitation.
\end{itemize}

% ══════════════════════════════════════════════════════════════════
\section{Methods}
% ══════════════════════════════════════════════════════════════════

\subsection{Notation}

Let $\mathbf{A} \in \mathbb{Z}_{\geq 0}^{m \times n}$ be a sparse count matrix with $m$ rows (genes), $n$ columns (cells), and $K = \text{nnz}(\mathbf{A})$ nonzero entries.
In CSC format, $\mathbf{A}$ is stored as three arrays: \texttt{indptr} (length $n+1$, int32), \texttt{indices} (length $K$, int32), and \texttt{data} (length $K$, float64).
The raw CSC size is:
\begin{equation}
S_{\text{CSC}} = 4(n+1) + 4K + 8K = 4(n+1) + 12K \;\text{bytes}
\label{eq:csc_size}
\end{equation}

\subsection{VOCSC encoding}

\subsubsection{Row reordering}
Rows (genes) are permuted in descending order of nonzero count $\text{nnz}_i = |\{j : A_{ij} > 0\}|$.
This concentrates the densest rows at small indices, reducing the magnitude of gaps in subsequent encoding.
The permutation $\pi$ is stored as a compressed side-channel.

\subsubsection{Value-ordered grouping}
For each column $j$, entries are grouped by value.
Let $v_1 < v_2 < \ldots < v_G$ be the $G_j$ distinct nonzero values in column $j$.
For each value $v_g$, the rows $\{i : A_{ij} = v_g\}$ (after permutation) form a sorted set $R_{j,g} = \{r_1 < r_2 < \ldots < r_{c_g}\}$.
The metadata per column encodes: $G_j$, and for each group, $(v_g, c_g)$, all as variable-length integers (varints).

\subsubsection{Gap-1 encoding}
For each group $R_{j,g}$, the sorted row indices are gap-1 encoded:
\begin{equation}
\delta_k = \begin{cases}
r_1, & k = 1 \\
r_k - r_{k-1} - 1, & k > 1
\end{cases}
\label{eq:gap1}
\end{equation}
This transforms a monotone sequence into a non-negative integer sequence centered near zero.
The key insight is that after row reordering, the densest rows cluster at the beginning, producing many small gaps in the dominant value group ($v=1$).

\subsubsection{Adaptive byte-plane splitting}
Raw gaps are stored as 32-bit unsigned integers.
We observe that in scRNA-seq data, gap values are universally bounded by the number of rows ($m < 2^{16}$ for typical gene counts), allowing a 16-bit representation.
The codec adaptively selects 2-plane or 4-plane byte splitting based on the observed maximum gap.

For 2-plane splitting ($\max(\delta) < 65536$), the $N$ gap values are reorganized as:
\begin{equation}
\text{Split}_2(\delta) = [\delta_1 \bmod 256, \ldots, \delta_N \bmod 256, \lfloor \delta_1/256 \rfloor, \ldots, \lfloor \delta_N/256 \rfloor]
\end{equation}
This groups bytes of similar magnitude, yielding highly compressible planes: the high byte is predominantly zero (94\% of gaps fit in a single byte), while the low byte inherits the entropy of the gap distribution.

\subsubsection{Chunked compression}
Columns are processed in chunks of $C = 1024$ (default).
Each chunk independently:
\begin{enumerate}
\item VOCSC-encodes its columns into metadata + gap arrays.
\item Byte-splits the gap array into 2 or 4 planes.
\item Concatenates metadata and split planes into a single buffer.
\item Compresses the buffer as one zstd frame at level 3.
\end{enumerate}
Chunking enables parallel decoding while providing sufficient context for the compressor.

\subsection{File format}

The on-disk format consists of:
\begin{itemize}
\item A 96-byte header containing matrix dimensions, chunk parameters, format flags, and blob sizes.
\item Compressed row permutation (zstd level~3).
\item Compressed column pointer deltas (zstd level~3).
\item Chunk size table ($4 \times N_c$ bytes).
\item $N_c$ compressed chunk blobs, each prefixed with 12 bytes of metadata (gap count, metadata size, compressed size).
\item An 8-byte footer (checksum placeholder + magic number).
\end{itemize}

\subsection{Decoding}

Decoding is fully parallelizable across chunks via OpenMP.
Each thread:
\begin{enumerate}
\item Decompresses its chunk blob (single zstd frame).
\item Separates the metadata prefix from the byte-split gap planes.
\item Reconstructs 32-bit gaps via byte-plane interleaving.
\item Walks the VOCSC metadata, scattering (row, value) pairs to the output CSC arrays with prefetching for the permutation lookup.
\end{enumerate}

\subsection{Complexity analysis}

\textbf{Encoding:} $O(K \log K)$ dominated by per-column sorting of (value, row) pairs.
VOCSC metadata and gap computation are $O(K)$.
Byte-split is $O(K)$.
Zstd compression is $O(K)$ with low constant.

\textbf{Decoding:} $O(K)$---a single pass over the gap stream with constant-time operations per element.
The byte-unsplit step is $O(K)$ with no branching; the VOCSC scatter is $O(K)$ with one table lookup (permutation) per gap.

% ══════════════════════════════════════════════════════════════════
\section{Experimental Setup}
% ══════════════════════════════════════════════════════════════════

\subsection{Corpus}

We evaluate on 20 scRNA-seq count matrices from the Gene Expression Omnibus (GEO)~\cite{edgar2002gene}, representing diverse tissues, species, and sequencing technologies.
Table~\ref{tab:corpus} summarizes the corpus.
Matrices range from 1.4M to 7.7M nonzeros, with density $0.17$--$5.9$\% and fraction of unit-valued entries $50$--$84$\%.

\begin{table}[t]
\centering
\caption{Corpus summary statistics ($n=20$ matrices).}
\label{tab:corpus}
\footnotesize
\begin{tabular}{lrrrr}
\toprule
Statistic & Median & Min & Max & Unit \\
\midrule
Rows (genes)        & 115,818 & 33,694 & 171,540 & --- \\
Columns (cells)     &   2,457 &    630 &  21,187 & --- \\
Nonzeros            &    3.6M &   1.4M &    7.7M & --- \\
Density             &  0.82\% & 0.17\% &  5.93\% & --- \\
Fraction $v{=}1$    &  76.5\% & 50.5\% &  84.1\% & --- \\
Fraction $v{\leq}3$ &  92.4\% & 78.5\% &  96.6\% & --- \\
Max value           &   3,058 &    368 &  29,933 & --- \\
Gap entropy         &    6.30 &   4.92 &    7.74 & bits \\
Raw CSC size        &  42.7 &  17.1 &  92.1 & MB \\
\bottomrule
\end{tabular}
\end{table}

\subsection{Baselines}

We compare VOCSC against general-purpose compressors applied to int32-cast CSC arrays (indptr + indices + data), which represents the practical floor for any CSC compression scheme:

\begin{itemize}
\item \textbf{gzip-6}~\cite{deutsch1996gzip}: Standard deflate, the most common general-purpose compressor.
\item \textbf{LZ4}~\cite{collet2013lz4}: Fast decompression-oriented compressor.
\item \textbf{zstd-3}~\cite{collet2018zstandard}: Fast Zstandard at the same level as our backend.
\item \textbf{zstd-22}~\cite{collet2018zstandard}: Maximum Zstandard compression level.
\item \textbf{Brotli-11}~\cite{alakuijala2019brotli}: Maximum Brotli compression quality.
\end{itemize}

All baselines compress the concatenated raw bytes of (indptr, indices, data) as a single stream.
No domain-specific preprocessing is applied.

\subsection{Hardware and measurement}

Experiments run on an Intel Xeon compute node with 256~GB RAM and 8 threads.
Encoding time is wall-clock; decoding time is the best of 10 iterations (measuring throughput under warm cache).
Throughput is reported as raw CSC bytes per second ($S_{\text{CSC}} / t$).

% ══════════════════════════════════════════════════════════════════
\section{Results}
% ══════════════════════════════════════════════════════════════════

\subsection{Compression performance}

Table~\ref{tab:main} compares VOCSC against baseline compressors across the 20-matrix corpus.
VOCSC achieves a median compression ratio of 12.5$\times$ over raw CSC (float64 values), with aggregate ratio of 12.9$\times$.
This represents a 2.2$\times$ improvement over zstd-3 applied to raw CSC (5.85$\times$) and a 1.4$\times$ improvement over the best generic compressor (Brotli-11 at 8.98$\times$, which requires 400$\times$ longer encoding).

\begin{table}[t]
\centering
\caption{Compression comparison on 20 scRNA-seq matrices. Ratios are vs.\ raw CSC (float64 values + int32 indices). Throughput at 8 threads.}
\label{tab:main}
\footnotesize
\begin{tabular}{lrrr}
\toprule
Method & Ratio & Enc.\ (MB/s) & Dec.\ (MB/s) \\
\midrule
Raw CSC (baseline)  &  1.00$\times$ & --- & --- \\
LZ4                 &  2.73$\times$ &  747 & --- \\
gzip-6              &  5.42$\times$ &   25 & --- \\
zstd-3 (raw CSC)    &  5.85$\times$ &  296 & --- \\
zstd-22 (raw CSC)   &  8.43$\times$ &    3 & --- \\
Brotli-11 (raw CSC) &  8.98$\times$ &    1 & --- \\
\midrule
\textbf{VOCSC + zstd-3} & \textbf{12.94}$\times$ & \textbf{240} & \textbf{4,624} \\
\bottomrule
\end{tabular}
\end{table}

Decoding throughput of 4,624~MB/s means a 43~MB matrix (3.6M nonzeros) decompresses in approximately 9~ms, enabling interactive analytical workflows.

\subsection{Ablation study}
\label{sec:ablation}

We decompose the VOCSC pipeline into stages and measure the cumulative compression gain at each step (Table~\ref{tab:ablation}).
The ablation is performed on 10 matrices from the corpus.

\begin{table}[t]
\centering
\caption{Ablation study: cumulative compression by pipeline stage. Sizes are summed across 10 matrices.}
\label{tab:ablation}
\footnotesize
\begin{tabular}{lrrr}
\toprule
Stage & Size (MB) & Ratio vs.\ CSC & Stage gain \\
\midrule
CSC (float64 values)     & 300.5 & 1.00$\times$ & --- \\
CSC (int32 values)       & 200.4 & 1.50$\times$ & 1.50$\times$ \\
VOCSC gap encoding       & 100.1 & 3.00$\times$ & 2.00$\times$ \\
+ Byte-plane split       &  50.1 & 6.00$\times$ & 2.00$\times$ \\
+ zstd-3                 &  20.8 & 14.4$\times$ & 2.41$\times$ \\
\midrule
\emph{Higher-effort variants:} & & & \\
\quad + zstd-22          &  20.4 & 14.7$\times$ & +2.1\% \\
\quad + Brotli-11        &  19.0 & 15.8$\times$ & +8.8\% \\
\bottomrule
\end{tabular}
\end{table}

Each stage contributes approximately a 2$\times$ size reduction.
The transition from CSC (float64) to CSC (int32) eliminates 4 bytes per nonzero from the value stream (values in count data are always small integers, never requiring 64-bit representation).
VOCSC gap encoding replaces the (index, value) pair with a single gap integer per nonzero, halving the per-element cost.
Byte-plane splitting reorganizes the gap bytes into homogeneous planes that compress 2.4$\times$ better under zstd-3 than unsplit gaps.

The marginal gains from higher-effort backends are modest: zstd-22 improves the ratio by 2.1\% over zstd-3 at the cost of 7$\times$ slower encoding, while Brotli-11 improves by 8.8\% at 30$\times$ slower encoding and 3$\times$ slower decoding.
We therefore recommend zstd-3 as the sole backend.

\subsection{Extended encoding variants}
\label{sec:extended}

We evaluate three refinements to the VOCSC encoding that exploit additional statistical structure in count matrices (Table~\ref{tab:extended}).

\textbf{Implicit value labels.}
In scRNA-seq data, the value groups within each column typically progress $1, 2, 3, \ldots$ (because 1 is the most common nonzero, 2 the second most common, etc.).
Rather than encoding each value label as a varint, we can assume consecutive progression from 1 and omit the label entirely.
When the sequence breaks, we switch to explicit (value, count) metadata.
Across 20 matrices, a median of 55.5\% of value groups are consecutive, but the metadata savings are small (varint 1--127 costs only a single byte per group).

\textbf{Dense head.}
After row density sorting, the first $D$ rows of each column are nearly fully populated: for $D = 128$, the head density exceeds 80\% on 15 of 20 matrices.
For these rows, we replace VOCSC gap encoding with a $D$-byte dense vector where each byte stores the cell value (0 for absent; values above 255 are clamped).
The remaining $m - D$ rows are encoded with standard VOCSC.
Dense encoding costs $D$ bytes per column regardless of occupancy, whereas VOCSC costs approximately 1.2 bytes per nonzero; the breakeven occurs at roughly 80\% head density.

\textbf{Hybrid.}
We combine all three ideas: a dense head of $D$ rows, implicit-value VOCSC for the consecutive groups in the sparse tail, and explicit VOCSC for any remaining non-consecutive groups.

\begin{table}[t]
\centering
\caption{Extended encoding variants vs.\ baseline VOCSC on 20 matrices. $\Delta$ is relative to baseline compressed size (lower is better). Dense/hybrid variants parameterized by head depth $D$.}
\label{tab:extended}
\footnotesize
\begin{tabular}{lrrrr}
\toprule
Variant & Median $\Delta$ & Mean $\Delta$ & Min $\Delta$ & Max $\Delta$ \\
\midrule
Implicit values    & $-$0.3\% & $-$0.3\% & $-$1.0\% & $-$0.2\% \\
Dense ($D{=}128$)  & $-$1.2\% & $-$1.4\% & $-$3.6\% & $-$0.1\% \\
Hybrid ($D{=}128$) & $-$1.3\% & $-$1.7\% & $-$3.9\% & $-$0.3\% \\
\midrule
Dense ($D{=}256$)  & $-$1.2\% & $-$0.9\% & $-$2.9\% & $+$3.1\% \\
Hybrid ($D{=}256$) & $-$1.4\% & $-$1.1\% & $-$3.1\% & $+$2.6\% \\
\bottomrule
\end{tabular}
\end{table}

The hybrid variant at $D = 128$ achieves the most consistent improvement: median $-$1.3\%, never worse than baseline (worst case $-$0.3\%).
At $D = 256$, the median improvement is larger ($-$1.4\%) but the scheme regresses on very sparse matrices (density $< 0.001$), where the 256-byte dense vector exceeds the VOCSC cost for partially-empty head rows.

Per-matrix, the best encoding always uses the hybrid scheme, with best-$D$ selection yielding improvements ranging from $-$0.6\% to $-$3.9\% (median $-$1.7\%).
However, the optimal $D$ depends on matrix statistics: dense matrices favor $D = 256$ while sparse matrices favor $D = 32$--$128$.

\textbf{Assessment.}
The extended variants provide a real but modest 1--4\% improvement over the baseline VOCSC pipeline.
This gain is substantially smaller than the 24\% improvement from byte-plane splitting or the 100\% from VOCSC grouping itself.
Moreover, adaptive $D$ selection adds codec complexity and would require per-chunk or per-matrix heuristics at encoding time.
We include these results to characterize the limits of structural exploitation in this domain, but do not adopt them in the reference implementation.
The current VOCSC pipeline already captures the dominant redundancy; these refinements operate in the diminishing-returns regime of the entropy curve.

\subsection{Why byte-plane splitting is essential}

Byte-plane splitting is the single most impactful preprocessing step after VOCSC encoding itself.
Table~\ref{tab:bsplit} shows the effect of applying the same backend compressor with and without byte-split preprocessing.

\begin{table}[t]
\centering
\caption{Impact of byte-plane splitting on compression ratio (10 matrices, summed). ``Raw'' = gaps compressed directly; ``Split'' = byte-split then compressed.}
\label{tab:bsplit}
\footnotesize
\begin{tabular}{lrrrr}
\toprule
Backend & Raw (MB) & Split (MB) & Raw ratio & Split ratio \\
\midrule
zstd-3   & 25.9 & 20.8 & 11.6$\times$ & 14.4$\times$ \\
zstd-22  & 23.7 & 20.4 & 12.7$\times$ & 14.7$\times$ \\
Brotli-11& 22.4 & 19.0 & 13.4$\times$ & 15.8$\times$ \\
\bottomrule
\end{tabular}
\end{table}

The improvement from byte-splitting is 24\% for zstd-3, 16\% for zstd-22, and 18\% for Brotli-11.
The mechanism is straightforward: gap values in scRNA-seq data are overwhelmingly small (94.2\% fit in a single byte).
Byte-split segregates the near-zero high byte-plane from the informative low byte-plane, allowing the compressor to exploit the extreme redundancy in the high plane separately.

\subsection{Compression level analysis}

Tables~\ref{tab:level} and~\ref{tab:brotli} show the Pareto frontier across zstd levels 1--22 and Brotli qualities 1--11.

\begin{table}[t]
\centering
\caption{Zstd level sweep (10 matrices, aggregate). Ratio is vs.\ raw CSC.}
\label{tab:level}
\footnotesize
\begin{tabular}{rrrc}
\toprule
Level & Ratio & Enc.\ (MB/s) & Dec.\ (MB/s) \\
\midrule
 1 & 12.17$\times$ & 211 & 4,125 \\
 3 & 12.08$\times$ & 198 & 3,959 \\
 7 & 12.18$\times$ & 164 & 3,599 \\
12 & 12.22$\times$ & 121 & 3,776 \\
15 & 12.23$\times$ & 90 & 3,711 \\
19 & 12.08$\times$ & 41 & 3,543 \\
22 & 12.09$\times$ & 27 & 3,591 \\
\bottomrule
\end{tabular}
\end{table}

\begin{table}[t]
\centering
\caption{Brotli quality sweep (10 matrices, aggregate).}
\label{tab:brotli}
\footnotesize
\begin{tabular}{rrrc}
\toprule
Quality & Ratio & Enc.\ (MB/s) & Dec.\ (MB/s) \\
\midrule
 1 & 12.13$\times$ &  155 & 1,659 \\
 5 & 12.18$\times$ &  130 & 1,756 \\
 9 & 12.20$\times$ &   58 & 1,752 \\
11 & 13.08$\times$ &    6 & 1,397 \\
\bottomrule
\end{tabular}
\end{table}

Across the entire zstd range (levels 1--22), the compression ratio varies by less than 1.3\% (12.08--12.23$\times$), while encoding speed drops from 211 to 27~MB/s ($7.8\times$ slower).
Brotli-11 achieves a meaningful improvement (13.08$\times$) but at 33$\times$ slower encoding and 2.8$\times$ slower decoding.
The flat ratio curve across zstd levels indicates that the VOCSC + byte-split preprocessing has already extracted most of the redundancy: the backend compressor merely exploits residual entropy with diminishing returns at higher levels.

\subsection{Synthetic data experiments}
\label{sec:synthetic}

We validate VOCSC behavior on synthetic sparse matrices where the statistical parameters are precisely controlled.

\subsubsection{Effect of density}

Table~\ref{tab:density} shows compression ratio as a function of matrix density (fraction of nonzero entries), with all other parameters held constant ($m = 30000$, $n = 5000$, negative binomial values with $r = 0.5$, $p = 0.3$).

\begin{table}[t]
\centering
\caption{Synthetic experiment: compression ratio vs.\ matrix density.}
\label{tab:density}
\footnotesize
\begin{tabular}{rrr}
\toprule
Density & nnz & VOCSC ratio \\
\midrule
0.001 &    150K & 5.14$\times$ \\
0.002 &    300K & 6.13$\times$ \\
0.005 &    748K & 7.44$\times$ \\
0.010 &  1,493K & 8.45$\times$ \\
0.020 &  2,970K & 9.59$\times$ \\
0.050 &  7,315K & 11.16$\times$ \\
0.100 & 14,274K & 12.55$\times$ \\
\bottomrule
\end{tabular}
\end{table}

Compression ratio increases with density following a logarithmic trend: $R \approx 3.29 \log_{10}(d) + 13.3$ ($R^2 = 0.998$).
Higher density means more nonzeros per column, producing more gaps per value group.
These gaps have smaller mean values (more tightly packed rows), yielding lower entropy and better compressibility.
At extremely low density (0.1\%), the overhead of metadata (group headers, permutation, pointers) becomes significant relative to the gap stream, reducing the effective ratio.

\subsubsection{Effect of value distribution}

\begin{table}[t]
\centering
\caption{Synthetic experiment: compression ratio vs.\ negative binomial $r$ parameter (controlling value dispersion).}
\label{tab:nb}
\footnotesize
\begin{tabular}{rrrr}
\toprule
$r$ & Frac.\ $v{=}1$ & Mean val & VOCSC ratio \\
\midrule
0.1 & 88.2\% &  1.2 &  9.91$\times$ \\
0.2 & 78.3\% &  1.5 &  9.36$\times$ \\
0.5 & 54.6\% &  2.2 &  8.45$\times$ \\
1.0 & 29.8\% &  3.3 &  7.86$\times$ \\
2.0 &  9.0\% &  5.7 &  7.35$\times$ \\
5.0 &  0.2\% & 12.7 &  6.98$\times$ \\
10.0 & 0.0\% & 24.4 &  6.51$\times$ \\
\bottomrule
\end{tabular}
\end{table}

When $r$ is small (high dispersion, many unit values), VOCSC achieves nearly 10$\times$ compression because the dominant $v{=}1$ group produces a single long gap stream with strong locality.
As $r$ increases, values spread across more groups, each with its own gap sequence, reducing per-group stream length and compressibility.
Even with no unit-valued entries ($r = 10$), VOCSC maintains 6.5$\times$ compression, demonstrating robustness beyond the scRNA-seq regime.

\subsubsection{Effect of matrix dimensions}

Compression ratio increases modestly with row count: from 7.12$\times$ at $m = 5000$ to 7.54$\times$ at $m = 100000$ (density = 0.5\%).
The mechanism is indirect: more rows at fixed density yield the same number of nonzeros but distributed across a larger row space, producing slightly larger gaps that compress marginally better under byte-split (more zero-valued high bytes).
Beyond 100K rows, the ratio plateaus.

\subsubsection{Effect of value distribution family}

\begin{table}[t]
\centering
\caption{Compression ratio by value distribution family (density = 1\%).}
\label{tab:dist}
\footnotesize
\begin{tabular}{lrr}
\toprule
Distribution & Distinct values & VOCSC ratio \\
\midrule
Geometric ($p{=}0.6$) & 18 & 8.91$\times$ \\
Neg.\ binomial ($r{=}0.5$, $p{=}0.3$) & 36 & 8.45$\times$ \\
Uniform ($[1, 100]$) & 203 & 5.93$\times$ \\
\bottomrule
\end{tabular}
\end{table}

The uniform distribution, which lacks value concentration, still achieves 5.93$\times$ compression.
This ``baseline'' ratio comes entirely from the gap encoding and byte-split pipeline operating on the position stream, independent of value structure.
Concentrated distributions (geometric, negative binomial) add 2--3$\times$ additional compression from the value-ordered grouping.

\subsection{Predictive compressibility model}
\label{sec:model}

We fit a linear model to predict the achieved VOCSC compression ratio from six observable matrix statistics:
\begin{equation}
\hat{R} = \beta_0 + \beta_1 \log_{10}(d) + \beta_2 f_1 + \beta_3 \log_2(\bar{\delta}) + \beta_4 H_\delta + \beta_5 \log_{10}(m) + \beta_6 \log_{10}(n)
\end{equation}
where $d$ = density, $f_1$ = fraction of unit-valued entries, $\bar{\delta}$ = mean gap, $H_\delta$ = gap entropy (bits), $m$ = rows, $n$ = columns.

\begin{table}[t]
\centering
\caption{Predictive model coefficients ($R^2 = 0.986$, $n = 20$ matrices).}
\label{tab:model}
\footnotesize
\begin{tabular}{lr}
\toprule
Feature & Coefficient \\
\midrule
Intercept                & $+18.29$ \\
$\log_{10}(\text{density})$ & $+3.55$ \\
Fraction $v{=}1$           & $+0.41$ \\
$\log_2(\text{mean gap})$  & $+0.72$ \\
Gap entropy (bits)          & $-2.52$ \\
$\log_{10}(m)$              & $+0.20$ \\
$\log_{10}(n)$              & $+3.36$ \\
\bottomrule
\end{tabular}
\end{table}

The model achieves $R^2 = 0.986$ on the 20-matrix corpus, confirming that compression ratio is well-determined by observable statistics.
The two strongest predictors are gap entropy ($\beta_4 = -2.52$, higher entropy $\Rightarrow$ lower ratio) and column count ($\beta_6 = +3.36$, more columns $\Rightarrow$ more chunks $\Rightarrow$ better compression).
Density enters with a positive coefficient ($+3.55$), consistent with the synthetic experiments: denser matrices produce more gaps per chunk, improving compressor context.

This model enables a priori estimation of file sizes and can guide storage allocation without trial compression.

% ══════════════════════════════════════════════════════════════════
\section{Discussion}
% ══════════════════════════════════════════════════════════════════

\subsection{Why VOCSC outperforms generic compressors}

The 2.2$\times$ advantage of VOCSC over zstd-3 on raw CSC arrays (12.9$\times$ vs.\ 5.85$\times$) arises from three structural transformations invisible to generic compressors:

\textbf{Value elimination.} By grouping entries by value, the explicit value stream is replaced by compact varint metadata (group value + count), saving 4--8 bytes per nonzero. For a matrix with $K$ nonzeros and median 27 distinct values per column, the metadata overhead is approximately $3 \times 27 = 81$ bytes per column vs.\ $8K/n$ bytes for the raw value stream.

\textbf{Gap-1 transform.} Converting absolute row indices to gap-1 deltas reduces the dynamic range by orders of magnitude. For a matrix with $m = 115{,}818$ rows and density 0.82\%, the mean gap is $\approx 122$---well within a single byte for most entries.

\textbf{Byte-plane separation.} Byte-split transforms structured 16/32-bit integers into homogeneous byte streams. The high byte-plane (mostly zeros, 94\% of gaps $< 256$) compresses to near zero, while the low byte-plane carries the remaining entropy. This is equivalent to column-major storage of a conceptual $N \times 2$ (or $N \times 4$) byte matrix, exploiting per-column redundancy.

\subsection{The case for a single compression tier}

Our ablation and level-sweep experiments reveal a crucial insight: after VOCSC + byte-split preprocessing, the residual entropy is so low that backend compressor effort has minimal impact on compression ratio.
Across zstd levels 1--22, the ratio varies by only 1.3\%.
Brotli-11 achieves a more substantial 8.8\% improvement, but at 33$\times$ slower encoding (6~MB/s vs.\ 198~MB/s) and 2.8$\times$ slower decoding (1,397~MB/s vs.\ 3,959~MB/s).

For a practical system compressing hundreds of terabytes of single-cell data, the encoding throughput cost is decisive: compressing 1~TB at 240~MB/s takes 70 minutes; at 6~MB/s it takes 46 hours.
The 8.8\% size reduction saves approximately 7~GB per 80~GB of raw CSC---far less than the operational cost of the additional compute time.

We therefore implement and recommend a single tier: VOCSC + byte-split + zstd-3.
Higher-effort variants (zstd-22, Brotli-11) remain theoretically available by adjusting the compression level parameter, but we consider the tradeoff unfavorable for all practical scenarios.

\subsection{Theoretical limits}

The gap entropy of our corpus ranges from 4.92 to 7.74 bits per gap.
For a matrix with $K$ nonzeros at median gap entropy $H = 6.30$ bits, the entropy floor is:
\begin{equation}
S_{\min} = \frac{K \times H}{8} \text{ bytes}
\end{equation}
For a typical matrix with $K = 3.6$M, this gives $S_{\min} \approx 2.84$~MB, or $15.5\times$ compression vs.\ 44.1~MB CSC.
Our achieved 12.5$\times$ (median) is 80\% of this theoretical limit, with the gap attributable to metadata overhead, permutation storage, and the sub-optimal (but fast) compression of zstd at level~3.

The 2.2$\times$ advantage over the best generic compressor applied to raw CSC (zstd-3 at 5.85$\times$) demonstrates that domain-aware preprocessing---not merely a better entropy coder---is the key to high compression of sparse count data.

\subsection{Limitations}

VOCSC is specialized for sparse integer count matrices.
It does not handle:
\begin{itemize}
\item Continuous-valued data (e.g., normalized expression).
\item Dense matrices (density $>$ 10--20\%).
\item Non-integer values.
\end{itemize}
For such data, the value-ordered grouping degenerates (every entry is a unique group), and the codec falls back to gap-only compression at approximately 6$\times$.

The current implementation requires loading the entire matrix into memory for encoding (due to row permutation computation).
Streaming encoding with approximate permutation is an area for future work.

% ══════════════════════════════════════════════════════════════════
\section{Related Work}
% ══════════════════════════════════════════════════════════════════

\textbf{General sparse matrix compression.}
The CSC and CSR formats~\cite{barrett1994templates} are the standard in-memory representations.
Libraries such as HDF5~\cite{folk2011overview} and Zarr~\cite{miles2024zarr} add chunk-level compression (gzip, LZ4, zstd) over raw array bytes.
SuiteSparse~\cite{davis2019algorithm} and MatrixMarket~\cite{boisvert1997matrix} define interchange formats but do not exploit data-level structure.

\textbf{Genomics formats.}
CRAM~\cite{hsi2011cram} and MPEG-G~\cite{voges2021mpeg} apply reference-based and column-specific compression to sequencing reads.
AnnData~\cite{anndata} uses HDF5 as a container with per-dataset compression.
CellDB~\cite{sikkema2023integrated} stores expression data in Parquet with run-length encoding.
None of these exploit the value-ordered structure of UMI count data.

\textbf{Integer compression.}
Byte-aligned encoding schemes such as VByte (varint)~\cite{williams1999compressing}, Group Varint~\cite{dean2009challenges}, and SIMD-BP128~\cite{lemire2015decoding} are widely used for inverted indices in information retrieval.
Byte-shuffle preprocessing for columnar integer data is used by Blosc~\cite{alted2010blosc} and bitshuffle~\cite{masui2015bitshuffle}, which transpose an $N \times w$ byte matrix (where $w$ is the word width) before compression.
Our byte-split step is a simplified variant tailored to the gap-1 output of VOCSC, operating at the granularity of entire compressed chunks rather than individual words.

\textbf{Single-cell data compression.}
ScZip~\cite{li2024sczip} applies delta coding and byte-split to CSC indices with zstd backend, achieving 5--7$\times$ compression.
The Cellarium project~\cite{cellarium2024} stores count matrices in a proprietary binary format with run-length encoding.
VOCSC extends beyond delta coding by eliminating the value stream entirely through value-ordered grouping, yielding the additional 2$\times$ factor.

% ══════════════════════════════════════════════════════════════════
\section{Conclusions}
% ══════════════════════════════════════════════════════════════════

We introduced VOCSC, a lossless compression scheme for sparse count matrices that achieves 13$\times$ compression over standard CSC at 4,624~MB/s decode throughput.
The codec exploits three domain-specific properties---value concentration, implicit values, and gap compressibility---through a simple four-stage pipeline.
Our ablation study shows that each stage contributes approximately 2$\times$, and that the preprocessing dominates: switching the backend compressor from zstd-3 to zstd-22 or Brotli-11 yields less than 9\% additional compression at 7--30$\times$ higher encoding cost.
A linear predictive model achieves $R^2 = 0.986$, enabling a priori file-size estimation.
Extended encoding variants---dense head rows, implicit value labels, and a hybrid scheme---yield a further median 1.3--1.7\% size reduction but operate in the diminishing-returns regime of the entropy curve and are not adopted in the reference implementation.
VOCSC is implemented as SinglePress, a header-only C++ library with Python bindings and OpenMP parallelism, available under the MIT license.
SinglePress uses the \texttt{.1pz} file extension and is designed for integration into production single-cell analysis pipelines.

% ── Acknowledgments ──────────────────────────────────────────────
\section*{Acknowledgments}
Compute resources provided by the Flatiron Institute.

% ── Data \& Code Availability ─────────────────────────────────────
\section*{Data \& Code Availability}
SinglePress is open-source under the MIT license.
Source code, Python bindings, and benchmark scripts are available at \url{https://github.com/Singlet-AI/singlepress}.
The \texttt{.1pz} file format specification and reference encoder/decoder are included in the repository.
Evaluation data are publicly available from the Gene Expression Omnibus (GEO).

% ── References ───────────────────────────────────────────────────
\bibliography{refs}

\end{document}
